%% Andrea Piccolo - Master CV
%% Comprehensive overview of all competencies, experience, and achievements.
%% Use as a master reference when tailoring role-specific CVs.
%%
%% Compile with: lualatex main_example.tex
%% (pdflatex often fails on modern MiKTeX with fontawesome5 font-expansion errors.)

\documentclass[11pt,a4paper,sans]{moderncv}
\moderncvstyle{banking}
\moderncvcolor{blue}

% Force both first and last name AND section headings to render in moderncv
% blue (color1). Default banking on lualatex+MiKTeX leaves these black, which
% looks inconsistent with the rest of the blue accent scheme.
\renewcommand*{\firstnamestyle}[1]{{\fontsize{34}{36}\bfseries\upshape\color{color1}#1}}
\renewcommand*{\lastnamestyle}[1]{{\fontsize{34}{36}\bfseries\upshape\color{color1}#1}}
\renewcommand*{\sectionstyle}[1]{{\sectionfont\color{color1}#1}}

\usepackage[utf8]{inputenc}
\usepackage{hyperref}
\hypersetup{
    colorlinks=true,
    linkcolor=blue,
    filecolor=magenta,
    urlcolor=blue,
    pdftitle={Andrea Piccolo - CV},
    pdfpagemode=FullScreen,
}
\usepackage[scale=0.77]{geometry}
\usepackage{needspace}
\usepackage{import}

% personal data
\name{Andrea}{Piccolo}
\address{Leeds, UK}{}{}
\phone[mobile]{+44 7831 007579}
\email{andreapiccolope@gmail.com}
\extrainfo{\href{https://www.linkedin.com/in/ap0889/}{LinkedIn}, \href{https://github.com/APiccolo89}{GitHub} \textbar{} UK Global Talent Visa -- no sponsorship required}

\begin{document}

\makecvtitle

% ============================================================
%     PROFILE STATEMENT
% ============================================================

\vspace{6pt}
\small{Python developer and computational scientist with a PhD in Geodynamics and 10 years of experience in large-scale simulation, data analysis, and automated scientific workflows. Skilled in processing TB-scale datasets, optimising numerical software (up to 10\texttimes{} speedup), and building reproducible, tested pipelines on HPC clusters. Looking to bring deep scientific rigour and robust Python development skills to applied data science or environmental modelling in industry.}

% ============================================================
%     CORE COMPETENCIES
% ============================================================

\section{Core Competencies}
\vspace{1pt}
\begin{itemize}

\item \textbf{Scientific Python Development}: NumPy, SciPy, pandas, h5py, Numba, pytest — numerical modelling, data analysis, visualisation, workflow automation

\item \textbf{Numerical Methods \& Simulation}: Finite Element Methods (FEniCSx), Finite Difference Methods, 2D/3D thermo-mechanical modelling, model optimisation and benchmarking

\item \textbf{Large-Scale Data Processing}: TB-scale simulation dataset management, statistical analysis, 2D/3D scientific visualisation, pipeline automation on HPC clusters (SLURM, PBS)

\item \textbf{Software Engineering Practices}: Git, GitHub Actions (CI/CD), regression testing, documentation (Sphinx, LaTeX, Markdown), reproducible workflows

\item \textbf{Earth Sciences Domain}: Computational geodynamics, crustal mechanics, basin modelling, geophysics — decade of research across four European institutions

\end{itemize}

% ============================================================
%     PROFESSIONAL EXPERIENCE
% ============================================================

\section{Professional Experience}
\vspace{3pt}
\begin{itemize}

\needspace{5\baselineskip}
\item{\cventry{02/2025--Present}{Research Fellow}{University of Leeds}{Leeds, UK}{}{\vspace{1pt}
\begin{itemize}
    \item Developed a numerical software package using FEniCSx and Python; optimised a pre-existing package achieving a 10\texttimes{} speedup
    \item Extended physical laws within the model to test novel geodynamic hypotheses
    \item Built automated workflows to run simulations, process outputs, and analyse large datasets
    \item Integrated regression tests to validate the numerical simulation against benchmarks
    \item Managed the full analytical workflow independently, from model design to documentation and reporting
\end{itemize}}}

\vspace{3pt}

\needspace{5\baselineskip}
\item{\cventry{11/2021--05/2024}{Research Fellow}{University of Bayreuth}{Bayreuth, Germany}{}{\vspace{1pt}
\begin{itemize}
    \item Delivered publicly available databases: 30 TB of simulated data condensed into 300 GB of accessible outputs
    \item Developed Python scripts for data cleaning, analysis, and visualisation to support project reporting
    \item Extended an open-source modelling package in Julia; merged forward mechanical modelling into basin forward modelling
    \item Collaborated across institutions, facilitating the KIT Karlsruhe basin modelling group with new simulation data
\end{itemize}}}

\vspace{3pt}

\needspace{5\baselineskip}
\item{\cventry{11/2015--07/2021}{Research Fellow}{Johannes Gutenberg Universität Mainz}{Mainz, Germany}{}{\vspace{1pt}
\begin{itemize}
    \item Designed a 0D reduced-order model to predict the behaviour of full 3D simulations, then built an optimisation algorithm using 2D experiment data to recover the parameter combinations reproducing that output
    \item Generated and processed large-scale simulation datasets (100+ TB data processed); explored and visualised data to identify physical trends
    \item Developed modelling scripts and data-processing routines, including 2D and 3D visualisation tools
    \item Incorporated chemical modelling into large-scale 3D numerical code
    \item Presented results at international conferences; supported supervision and teaching activities
\end{itemize}}}

\end{itemize}

% ============================================================
%     EDUCATION
% ============================================================

\section{Education}
\vspace{1pt}
\begin{itemize}

\needspace{5\baselineskip}
\item{\cventry{11/2015--05/2019}{PhD in Computational Geodynamics}{Johannes Gutenberg Universität Mainz}{Mainz, Germany}{}{\vspace{1pt}
Thesis: ``Continental crust production during the early Earth: Constraints from geodynamic models.''
}}

\vspace{3pt}

\item{\cventry{10/2012--03/2015}{MSc in Geology and Technical Geology}{University of Padua}{Padua, Italy}{}{\vspace{1pt}
Data analysis (chemical, petrological, stratigraphical), numerical modelling, scientific communication, field data acquisition.
}}

\vspace{3pt}

\item{\cventry{10/2008--03/2012}{BSc in Earth Sciences}{University of Padua}{Padua, Italy}{}{\vspace{1pt}
Data analysis, MATLAB, field data acquisition.
}}

\end{itemize}

% ============================================================
%     LANGUAGES
% ============================================================

\section{Languages}
\vspace{1pt}
\begin{itemize}
\item Italian (native), English (C1), German (B1).
\end{itemize}

% ============================================================
%     REFERENCES
% ============================================================

\section{References}
\vspace{1pt}
\begin{itemize}
\item Prof.\ Dr.\ Timothy J.\ Craig --- University of Leeds (\href{mailto:T.J.Craig@leeds.ac.uk}{T.J.Craig@leeds.ac.uk})
\item Prof.\ Dr.\ Marcell Thielmann --- University of Bonn (\href{mailto:thielmann@geo.uni-bonn.de}{thielmann@geo.uni-bonn.de})
\item Prof.\ Dr.\ Nevena Andric Tomasevic --- KIT Karlsruhe (\href{mailto:nevena.tomasevic@kit.edu}{nevena.tomasevic@kit.edu})
\item Dr.\ Arne Spang --- University of Bayreuth (\href{mailto:arne.spang@uni-bayreuth.de}{arne.spang@uni-bayreuth.de})
\end{itemize}

\end{document}
